
Persistenz

Damit sind alle wesentlichen Informationen bekannt, die man benötigt,
um mit Daten sinnvoll arbeiten zu können.

Bisher wurden jedoch nur flüchtige (transiente) Daten zur Laufzeit
(runtime) betrachtet. Will man sie dauerhaft (persistent) speichern,
z. B. vom Arbeitsspeicher (RAM) in eine Datei auf Festplatte,
oder in umgekehrter Richtung einlesen, so sind einige zusätzliche
Informationen nötig:

channel
encoding
language (file format)
format (value format)

Die Tab. xx fasst alle neun in diesem Artikel besprochenen Merkmale
zur Beschreibung von Daten zusammen.

TABELLE

Bezeichnung         Beschreibung/Synonyme                           Beispiel
name                Name, Bezeichner (variable), Schlüssel (key)
channel             Kommunikationskanal                             inline, file
encoding            Kodierung                                       utf-8
language            Sprache, Dateiformat                            text/plain, text/cybol
format              Wertformat, MIME type                           text/plain, number/integer
type                Typ, Art der Abstraktion                        (nur intern verwendet)
model               Attribut, Eigenschaft, Member, Field, Wert (value)
properties          "
reference           Garbage Collection
